\chapter{Methodology}

This chapter explains how the project was conducted and why the selected methodological choices were appropriate for the research context.
The methodological design combines iterative software development, stakeholder driven refinement, and controlled validation in a Proof of Concept (PoC) setting.

The chapter focuses on how the work was performed, how decisions were justified, and how evidence was collected.
Detailed architecture and implementation are presented in later chapters, while this chapter establishes the methodological logic behind those outcomes.

\section{Development Approach (Agile and Iterative Development)}

The project followed an Agile and iterative development approach, as requirements evolved during the project period through continuous feedback from stakeholders and the academic supervisor \cite{agile_manifesto}.
This choice was not only practical but methodological: AI assisted systems involve technical uncertainty, where output quality and integration behavior are not fully predictable at project start.

A plan-driven approach such as Waterfall was considered less suitable.
Waterfall assumes stable, fully specified requirements and low discovery during execution, while this project required repeated adjustments related to workflow governance, retrieval behavior, model reliability, and operational constraints.

\vspace{0.4cm}

An MVP strategy was used to establish a functional baseline early.
After baseline stability was reached, additional functionality was introduced incrementally in short cycles.
Each cycle included planning, implementation, verification, and review.
This structure enabled early detection of integration issues and reduced the risk of late-stage redesign.

The iterative process also improved traceability.
Design and code changes were linked to specific requirements and observed behavior, making it possible to justify why features were added, modified, or deferred during the PoC period.

\begin{figure}[H]
\centering
\includegraphics[width=0.75\linewidth]{Figures/Agile_Methodology1.png}
\caption{Agile development lifecycle illustrating iterative and cyclic process stages \cite{GeeksforGeeks2025Jul}.}
\label{fig:agile_model}
\end{figure}

As illustrated in Figure~\ref{fig:agile_model}, the Agile lifecycle follows a cyclic process consisting of requirement definition, design, implementation, testing, and review phases\cite{agile_manifesto}.
In this project, each cycle produced a verifiable functional increment, which supported continuous alignment between project goals and implemented system behavior.

\section{Requirements Engineering and Prioritization Strategy}

Requirements engineering was treated as a methodological activity, not only a documentation step.
Before detailed design and coding, the project established a requirement baseline to define scope boundaries, quality expectations, and validation criteria for the PoC.
The baseline was then refined iteratively as implementation feedback exposed practical constraints related to runtime behavior, data quality, and integration complexity.

\vspace{0.4cm}

To keep requirements testable and traceable, each requirement was assigned a unique identifier and grouped into two categories:
\textbf{functional requirements} (what the system must do) and \textbf{non-functional requirements} (how reliably and safely the system must operate).
This distinction was methodologically important because the project objective was not only to implement features, but to demonstrate governed and dependable workflow behavior.

\vspace{0.4cm}

Prioritization was managed using a MoSCoW strategy in each iteration.
Requirements that directly enabled an end-to-end governed workflow (upload, review, approval, apply, retrievable knowledge) were treated as \textit{Must have}.
Requirements improving robustness and observability (for example fallback retrieval and indexing visibility) were treated as \textit{Should have}.
Lower-impact extensions were classified as \textit{Could have}, while out-of-scope enterprise-scale items were documented as \textit{Won't have (for now)}.

\begin{figure}[H]
\centering
\includegraphics[width=0.9\linewidth]{Figures/MoSCoW prioritization.png}
\caption{MoSCoW prioritization model showing classification of requirements based on importance \cite{KECG2023MoSCoW}.}
\label{fig:moscow}
\end{figure}

As shown in Figure~\ref{fig:moscow}, this method provided explicit scope control under limited thesis time and resource constraints.
Methodologically, it improved decision transparency because implementation order could be justified by requirement criticality instead of ad hoc coding preference.

\section{System Approach and Processing Strategy}

The system was developed using a staged processing strategy in which data passes through explicitly defined transformation and validation stages.
Methodologically, this was selected to improve modularity, reproducibility, and control of state transitions.

Rather than treating the platform as one monolithic process, the workflow was divided into bounded stages with explicit responsibilities and interfaces.
This aligns with established pipeline oriented system design principles, where separation of concerns enables clearer testing boundaries and improves experimental control.

\vspace{0.4cm}

This strategy provides three methodological advantages.
\textbf{First}, each stage can be validated independently before full integration.
\textbf{Second}, progression between stages can be governed through explicit rules, supporting quality control and auditability.
\textbf{Third}, failures in one stage can be isolated and handled without invalidating the entire workflow.

From a research perspective, staged processing also improves reproducibility: similar inputs can be traced through the same defined sequence of steps, making outcomes easier to compare and interpret across iterations.

\section{Human in the Loop Design Principle}

A human in the loop principle was adopted as a central methodological decision.
The AI component was used to generate structured suggestions, while final acceptance remained a human responsibility.

This choice was motivated by domain requirements: in industrial knowledge environments, correctness, accountability, and trust are more important than full automation speed.
The method therefore prioritizes assisted decision-making rather than autonomous decision-making.

\vspace{0.4cm}

Human review functions as a validation gate between AI generation and publication, reducing the risk of propagating incorrect or unverifiable content.
This is consistent with governance-oriented and trustworthy AI practices, where human oversight is recommended to mitigate risks such as hallucination, weak grounding, and contextually inappropriate outputs \cite{EU_HLEG_2019}.

In methodological terms, model output is treated as advisory evidence, while domain experts remain accountable for final state changes.

\section{Tool Selection and Technical Rationale}

Tool selection was guided by methodological and practical constraints: local execution, cost control, transparent data handling, reproducibility, and integration feasibility.
The selected stack was intended to balance experimental flexibility with engineering realism.

FastAPI was selected as backend framework because it supports modular API design, typed request validation, and clear service boundaries, which are beneficial for iterative development and controlled testing \cite{FastAPI}.
A web frontend was used to support interactive review and retrieval in one user workflow.
Workflow metadata and state transitions were persisted in a relational store, while semantic retrieval required vector indexing infrastructure \cite{Schwaber-Cohen2024Apr}.

\vspace{0.4cm}

A local model runtime was prioritized over commercial cloud APIs to reduce dependency on paid inference and to maintain control over sensitive processing contexts.
This choice also improves reproducibility because runtime behavior is less exposed to external API changes, pricing constraints, and remote-service variability.

Alternative configurations (larger hosted models, external APIs, and fully managed cloud stacks) were evaluated but not selected for the PoC due to trade-offs in cost, latency, operational complexity, and controllability.

\section{Project Management and Workflow (Kanban)}

Project execution was managed through a Kanban workflow in GitHub Projects \cite{kanban}.
Tasks were represented as cards and moved across defined states (for example backlog, ready, in progress, review, and done) \cite{our_Kanban_workflow_in_GitHub}.
This provided transparent progress tracking and dependency visibility.

Kanban was preferred over sprint-based Scrum because task duration and priority changed dynamically due to research driven discovery and integration uncertainty.
A continuous flow model was therefore more suitable than fixed length sprint commitments.

\vspace{0.3cm}

Kanban complemented iterative development by limiting work in progress and emphasizing completion before expansion.
Requirement priority was operationalized through MoSCoW-inspired planning, where core requirements were implemented first and lower-priority features were deferred when necessary \cite{productplan_moscow}.
Methodologically, this reduced scope drift and strengthened delivery discipline in the PoC setting.

\section{Methodological Strengths and Limitations}

The selected methodology has clear strengths.
It supports adaptability under evolving requirements, enables early delivery of a functional baseline, and promotes continuous validation through short feedback loops.
The staged processing strategy improves modular testing and traceability, while human in the loop review strengthens governance and trust.

\vspace{0.3cm}

The methodology also has limitations.
Iterative development can introduce temporary architectural inconsistency when priorities shift rapidly.
Local model execution constrains performance and may increase latency under heavier workloads.
The evaluation relies partly on expert judgment, which introduces subjectivity.
In addition, dataset diversity and project duration are limited, reducing generalizability beyond the PoC context.

\vspace{0.3cm}

Despite these limitations, the chosen method is considered appropriate for the thesis objective: developing and validating a robust, traceable, and practically relevant PoC under realistic academic and engineering constraints.

\section{Evaluation Strategy}

To evaluate the effectiveness of the proposed system, a qualitative and scenario-based evaluation strategy was applied.
The purpose was to systematically assess whether the system supports reliable document processing, controlled knowledge updates, and grounded information retrieval in realistic use conditions.

A qualitative strategy was selected because core success criteria include human interpretive judgment, governance behavior, and workflow correctness, which cannot be fully captured through isolated quantitative metrics.

\vspace{0.4cm}

This approach supports internal validity by ensuring that observed system behavior can be directly linked to predefined scenarios and explicit evaluation criteria.

\vspace{0.4cm}

Validation was implemented through a layered strategy combining focused automated tests and scenario-based API verification.
At the first level, targeted backend tests validated isolated behaviors such as authorization rules, apply consistency, and retrieval fallback logic.
At the second level, endpoint interactions were verified through realistic request-response flows with state transition checks.
At the third level, end-to-end scenarios were executed from document upload to approved publication and searchable retrieval.

\vspace{0.4cm}

This layered setup was selected because correctness depends on both component-level logic and cross-component orchestration.
In particular, requirements related to traceability, reliability, and observability can only be verified when workflow state, background indexing, and retrieval behavior are evaluated together.

The evaluation addressed three dimensions:

\vspace{0.4cm}
\textbf{1) Suggestion quality.}

Generated suggestions were reviewed by an expert user and compared against source documents.
Assessment criteria included relevance (coverage of key source concepts), factual correctness, structural clarity, and usefulness for approval decisions.

\textbf{2) Workflow correctness and traceability.}

End-to-end scenarios verified that only authorized users could approve/apply changes, that workflow states transitioned correctly, and that approved changes were persisted and auditable.

\textbf{3) Retrieval quality and runtime robustness.}

Scenarios verified that answers were grounded in approved knowledge, included source references, and remained operational under practical conditions such as indexing delays, timeout cases, and fallback activation.

\vspace{0.5cm}

Evaluation evidence was collected through manual inspection, structured scenario execution, endpoint level verification, and continuous observation during development iterations.
Results are presented in Chapter 4 together with requirement-traceable interpretation.

The evaluation has known limitations: it is not based on large benchmark datasets, and several judgments are expert dependent.
However, for a governance-oriented PoC, this strategy provides relevant and credible evidence of whether the integrated system behavior meets the intended objectives.

External validity is limited, as results are derived from a constrained PoC environment and may not fully generalize to large scale production settings.